% 乘积法则（综述）
% license CCBYSA3
% type Wiki

本文根据 CC-BY-SA 协议转载翻译自维基百科\href{https://en.wikipedia.org/wiki/Product_rule}{相关文章}

在微积分中，乘积法则（或称莱布尼茨法则[2]、莱布尼茨乘积法则）是一个用于求两个或多个函数乘积的导数的公式。对于两个函数，可以用拉格朗日记号表示为：
$$
(u \cdot v)' = u' \cdot v + u \cdot v',~
$$
或用莱布尼茨记号表示为：
$$
\frac{d}{dx}(u \cdot v) = \frac{du}{dx} \cdot v + u \cdot \frac{dv}{dx}.~
$$
该法则可以推广到三个或更多函数的乘积，推广到乘积的高阶导数规则，以及其他语境中。
\subsection{发现}
这一法则的发现通常归功于 戈特弗里德·莱布尼茨，他使用“无穷小”（现代微分的前身）进行了论证。[3]（然而，莱布尼茨论文的译者 J. M. Child[4] 认为应归功于艾萨克·巴罗。）以下是莱布尼茨的推理：[5]设$u$和$v$是函数。则$d(uv)$等于两个相继$uv$之差；其中一个是$uv$，另一个是$(u+du)(v+dv)$。于是：
$$
d(u \cdot v) = (u+du)\cdot(v+dv) - u\cdot v 
= u \cdot dv + v \cdot du + du \cdot dv.~
$$
由于项$du \cdot dv$相较于$du$和$dv$是“可忽略的”，莱布尼茨得出结论：
$$
d(u \cdot v) = v \cdot du + u \cdot dv,~
$$
这确实就是乘积法则的微分形式。

如果两边同时除以微分$dx$，就得到：
$$
\frac{d}{dx}(u \cdot v) = v \cdot \frac{du}{dx} + u \cdot \frac{dv}{dx},~
$$
这同样可以用拉格朗日记号写作：
$$
(u \cdot v)' = v \cdot u' + u \cdot v'.~
$$
\subsubsection{最初的证明}
莱布尼茨和牛顿都给出了乘积法则的证明，但以现代标准来看并不严格。
莱布尼茨借助“无穷小量”进行推理，把积解释为矩形的面积；而牛顿则使用“流动量”的概念进行推理。[6][7]
\subsection{例子}
\begin{itemize}
\item 假设我们要求函数$f(x) = x^{2}\sin(x)$的导数。应用乘积法则，可以得到$f'(x) = 2x \cdot \sin(x) + x^{2}\cos(x)$，因为$x^{2}$的导数是 $2x$，而正弦函数的导数是余弦函数。
\item 乘积法则的一个特殊情形是常数倍法则：如果 $c$ 是一个常数，而$f(x)$是一个可微函数，那么$(cf)(x) = c \cdot f(x)$也是可微的，其导数为$(cf)'(x) = c \cdot f'(x)$。这一结果可以直接从乘积法则推出，因为任何常数的导数都是零。
\item 这一点结合求导的和法则，表明微分运算是线性的。此外，分部积分公式就是由乘积法则推导出来的，商法则（弱形式）也是如此。所谓“弱形式”，指的是它并没有证明商一定可微，而只是指出：如果可微，其导数是什么。
\end{itemize}
\subsection{证明}
\subsubsection{导数的极限定义}
设$h(x) = f(x)g(x)$，并且假设$f$和$g$在$x$处可导。我们要证明$h$在 $x$ 处也可导，且其导数为$h'(x) = f'(x)g(x) + f(x)g'(x)$.为此，在分子中加上并减去一项$f(x)g(x+\Delta x) - f(x)g(x+\Delta x)$（其值为零，因此不改变原式），从而可以因式分解，然后利用极限的性质来处理。
$$
\begin{aligned}
h'(x) &= \lim_{\Delta x \to 0} \frac{h(x+\Delta x) - h(x)}{\Delta x} \\[6pt]
&= \lim_{\Delta x \to 0} \frac{f(x+\Delta x)g(x+\Delta x) - f(x)g(x)}{\Delta x} \\[6pt]
&= \lim_{\Delta x \to 0} \frac{f(x+\Delta x)g(x+\Delta x) - f(x)g(x+\Delta x) + f(x)g(x+\Delta x) - f(x)g(x)}{\Delta x} \\[6pt]
&= \lim_{\Delta x \to 0} \frac{\big(f(x+\Delta x)-f(x)\big)\cdot g(x+\Delta x) + f(x)\cdot \big(g(x+\Delta x)-g(x)\big)}{\Delta x} \\[6pt]
&= \lim_{\Delta x \to 0} \frac{f(x+\Delta x)-f(x)}{\Delta x} \cdot \lim_{\Delta x \to 0} g(x+\Delta x) \\ 
&\quad + \lim_{\Delta x \to 0} f(x) \cdot \lim_{\Delta x \to 0} \frac{g(x+\Delta x)-g(x)}{\Delta x} \\[6pt]
&= f'(x)g(x) + f(x)g'(x).
\end{aligned}~
$$
其中$\lim_{\Delta x \to 0} g(x+\Delta x) = g(x)$是因为可导函数必定连续。要不要我也帮你翻译利用微分的方法证明乘积法则的版本？
